%=======================================================================
%  IE 5311-001 (D01) Principles of Optimization
%  Homework 02 solutions
%
%  Packages used: geometry, amsmath, amssymb, amsthm, booktabs, array.
%  All ship with Overleaf and with a full TeX Live install.
%=======================================================================
\documentclass[11pt]{article}

\usepackage[margin=1in]{geometry}
\usepackage{amsmath}
\usepackage{amssymb}
\usepackage{amsthm}
\usepackage{booktabs}
\usepackage{array}

\DeclareMathOperator*{\minimize}{min}
\DeclareMathOperator*{\maximize}{max}
\DeclareMathOperator{\CVaR}{CVaR}
\DeclareMathOperator{\VaR}{VaR}
\DeclareMathOperator{\supp}{supp}
\DeclareMathOperator{\conv}{conv}
\DeclareMathOperator{\ext}{ext}

\newtheorem{lemma}{Lemma}
\newtheorem{proposition}{Proposition}
\theoremstyle{definition}
\newtheorem{remark}{Remark}

\setlength{\parskip}{0.45em}
\setlength{\parindent}{0pt}

\begin{document}

%=======================================================================
\begin{center}
  {\large IE 5311-001 (D01) Principles of Optimization}\\[4pt]
  Lidivine Kengne\\[2pt]
  Homework 02
\end{center}
\vspace{0.6em}

%=======================================================================
\section*{0.\quad Common model}
%=======================================================================

Problems 1 to 3 share one setting, so it is stated once here and reused.

\subsection*{0.1\quad Sets}

\begin{center}
\begin{tabular}{@{}ll@{}}
\toprule
Symbol & Definition \\
\midrule
$N = \{1,\dots,n\}$ & Products, indexed by $i$ \\
$\Xi \subseteq \mathbb{R}^n_+$ & Support of the demand vector, a polytope,
   $\Xi = \supp(P)$ \\
$\ext(\Xi) = \{v^1,\dots,v^K\}$ & Extreme points of $\Xi$, indexed by $k$ \\
$M = \{1,\dots,m\}$ & Sample indices in the SAA, indexed by $j$ \\
$S = \{1,\dots,|S|\}$ & Scenario indices when $\Xi$ is finite, indexed by $s$ \\
\bottomrule
\end{tabular}
\end{center}

\subsection*{0.2\quad Parameters}

\begin{center}
\begin{tabular}{@{}lll@{}}
\toprule
Symbol & Definition & Units \\
\midrule
$h_i \ge 0$ & Holding cost per leftover unit of product $i$ & \$/unit \\
$b_i \ge 0$ & Backlog cost per unsatisfied unit of product $i$ & \$/unit \\
$c_i = 0$   & Purchasing cost, assumed zero throughout & \$/unit \\
$\rho \in (0,1)$ & Half-width of the order-to-demand band & dimensionless \\
$\beta \in (0,1)$ & Confidence level of the risk measure & dimensionless \\
$p_s > 0$ & Probability of scenario $s$, $\sum_{s} p_s = 1$ & dimensionless \\
\bottomrule
\end{tabular}
\end{center}

\subsection*{0.3\quad Random element and decision variables}

\[
  \xi \in \Xi \subseteq \mathbb{R}^n_+ ,\quad \xi \sim P
  \qquad\text{(demand, units; observed after $x$ is chosen)}
\]
\[
  x_i \ge 0 \;=\; \text{order quantity of product } i \quad [\text{units}],
  \qquad i \in N .
\]

Auxiliary variables are introduced only where a reformulation needs them,
and are named at that point.

\subsection*{0.4\quad Cost function}

For an order $x$ and a realised demand $\xi$,
\begin{equation}
  f(x,\xi) \;=\; \sum_{i \in N}
    \Big[\, h_i\,(x_i - \xi_i)^{+} \;+\; b_i\,(\xi_i - x_i)^{+} \,\Big],
  \qquad (a)^{+} := \max\{a,0\}.
  \label{eq:cost}
\end{equation}
Since $c = 0$, no purchasing term appears. In compact vector form, with
$\odot$ the componentwise product and $\mathbf{1}$ the all-ones vector,
\[
  f(x,\xi) \;=\; \mathbf{1}^{\!\top}
    \Big[\, h \odot (x-\xi)^{+} + b \odot (\xi-x)^{+} \Big].
\]

Two structural facts are used repeatedly and are recorded now.

\begin{lemma}[Pointwise identity]
\label{lem:identity}
For any $h,b \ge 0$ and any $u \in \mathbb{R}$,
\[
  h\,(u)^{+} + b\,(-u)^{+} \;=\; \max\{\,h u,\; -b u\,\}.
\]
\end{lemma}

\begin{proof}
If $u \ge 0$ then the left side is $hu$, while $hu \ge 0 \ge -bu$, so the
right side is also $hu$. If $u < 0$ then the left side is $-bu$, while
$-bu \ge 0 \ge hu$, so the right side is also $-bu$. The two cases are
exhaustive.
\end{proof}

Applying Lemma~\ref{lem:identity} coordinatewise with $u = x_i - \xi_i$,
\begin{equation}
  f(x,\xi) \;=\; \sum_{i \in N}
    \max\big\{\, h_i (x_i - \xi_i),\; b_i(\xi_i - x_i) \,\big\}.
  \label{eq:costmax}
\end{equation}
This is the form every linearisation in this homework starts from. Note
also that each summand is nonnegative, so $f(x,\xi) \ge 0$ always.

\begin{lemma}[Convexity in each argument]
\label{lem:convex}
$f(\cdot,\xi)$ is convex in $x$ for fixed $\xi$, and $f(x,\cdot)$ is
convex in $\xi$ for fixed $x$. Both are piecewise linear.
\end{lemma}

\begin{proof}
By \eqref{eq:costmax} each summand is the maximum of two affine
functions, of $x$ for fixed $\xi$ and of $\xi$ for fixed $x$. A maximum
of finitely many affine functions is convex and piecewise linear, and a
sum of convex functions is convex.
\end{proof}

Convexity in $\xi$ is what makes Problem 2 finite. Convexity in $x$ is
what makes every model below a convex program.

%=======================================================================
\newpage
\section*{Problem 1.\quad Stochastic Multidimensional Newsvendor
  \quad[25 pts]}
%=======================================================================

\subsection*{1.1\quad(Part 1) The stochastic model}

No reformulation is performed; the expectation and the positive parts are
left as they stand.
\begin{align}
  \minimize_{x} \quad & \mathbb{E}_{P}\big[\, f(x,\xi) \,\big]
    \;=\; \mathbb{E}_{P}\!\left[\, \sum_{i \in N}
      \Big( h_i (x_i - \xi_i)^{+} + b_i (\xi_i - x_i)^{+} \Big) \right] \\
  \text{s.t.}\quad & x_i \ge 0 \qquad \forall\, i \in N .
\end{align}

This is a single-stage problem: $x$ is chosen before $\xi$ is observed and
there is no recourse decision. The randomness enters the objective only.

\subsection*{1.2\quad(Part 2) Sample average approximation}

Draw $\xi^1,\dots,\xi^m$ independently from $P$ and replace the
expectation by the empirical mean over the sample:
\begin{align}
  \minimize_{x \ge 0} \quad
    & \frac{1}{m} \sum_{j \in M} f(x,\xi^{j})
    \;=\; \frac{1}{m} \sum_{j \in M} \sum_{i \in N}
      \max\big\{ h_i (x_i - \xi^{j}_i),\; b_i(\xi^{j}_i - x_i) \big\}.
\end{align}

Using \eqref{eq:costmax}, introduce $t^{j}_i \in \mathbb{R}$ for the inner
maximum of product $i$ in sample $j$. The SAA problem is the linear
program
\begin{align}
  \minimize_{x,\,t} \quad & \frac{1}{m} \sum_{j \in M} \sum_{i \in N}
    t^{j}_i \\
  \text{s.t.}\quad
    & t^{j}_i \;\ge\; h_i\,(x_i - \xi^{j}_i)
      && \forall\, i \in N,\; j \in M \\
    & t^{j}_i \;\ge\; b_i\,(\xi^{j}_i - x_i)
      && \forall\, i \in N,\; j \in M \\
    & x_i \ge 0 && \forall\, i \in N .
\end{align}
Size: $n + nm$ variables and $2nm$ structural rows. The linearisation is
exact because the objective is being minimised and each $t^{j}_i$ is
pushed down to the larger of its two lower bounds.

\begin{remark}
The SAA optimal value is the \emph{in-sample} average at the \emph{same}
sample that defined it, so it is optimistically biased downward. The
honest out-of-sample number is $\mathbb{E}_P[f(\hat{x}_m,\xi)]$ evaluated
at the SAA solution $\hat{x}_m$. Section~1.7 reports both.
\end{remark}

\subsection*{1.3\quad(Part 3a) Physical meaning of the band constraint}

The additional constraint is
\begin{equation}
  (1-\rho)\,\xi_i \;\le\; x_i \;\le\; (1+\rho)\,\xi_i
  \qquad \forall\, i \in N .
  \label{eq:band}
\end{equation}

Read one product at a time. The upper bound says the order may exceed
demand by at most a fraction $\rho$: it caps deliberate over-stocking, and
therefore caps waste and the capital tied up in leftovers. The lower bound
says the order may fall short of demand by at most a fraction $\rho$: it
is a service-level floor, guaranteeing that at least a fraction $1-\rho$
of demand is covered from stock. Together they force the order to track
demand within a relative tolerance band of half-width $\rho$, so $\rho$ is
a dimensionless mismatch tolerance. At $\rho \to 0$ the order must equal
demand exactly; at $\rho \to 1$ the lower bound disappears and the upper
bound allows ordering twice the demand.

The important subtlety is one of timing. The vector $x$ is a here-and-now
decision, fixed before $\xi$ is seen, while \eqref{eq:band} refers to the
realised $\xi$. A single deterministic $x$ cannot adapt to the
realisation, so the constraint can only be imposed in a probabilistic
sense. That is what Part 3b makes precise and what Part 3c then collapses.

\subsection*{1.4\quad(Part 3b) The constrained stochastic model}

Requiring the band to hold almost surely under $P$,
\begin{align}
  \minimize_{x} \quad & \mathbb{E}_{P}\big[ f(x,\xi) \big] \\
  \text{s.t.}\quad
    & (1-\rho)\,\xi_i \le x_i \le (1+\rho)\,\xi_i
      \quad \forall\, i \in N, \qquad P\text{-a.e.} \\
    & x_i \ge 0 \qquad \forall\, i \in N ,
\end{align}
where the probability constraint means
\[
  P\Big( \big\{ \xi \in \Xi :\;
    (1-\rho)\xi_i \le x_i \le (1+\rho)\xi_i \;\;\forall i \in N \big\} \Big)
  \;=\; 1 .
\]

\subsection*{1.5\quad(Part 3c) The a.e.\ constraint is a robust constraint}

\begin{proposition}
\label{prop:aerobust}
Let $\Xi = \supp(P)$ be compact and let $x \in \mathbb{R}^n_+$ be fixed.
The following are equivalent.
\begin{enumerate}
  \item[(i)] \eqref{eq:band} holds $P$-a.e.
  \item[(ii)] \eqref{eq:band} holds for every $\xi \in \Xi$.
  \item[(iii)] For every $i \in N$,
    \[
      x_i \;\le\; (1+\rho)\,\min_{\xi \in \Xi} \xi_i
      \qquad\text{and}\qquad
      x_i \;\ge\; (1-\rho)\,\max_{\xi \in \Xi} \xi_i .
    \]
\end{enumerate}
\end{proposition}

\begin{proof}
\textbf{(ii) $\Rightarrow$ (i).} Since $\Xi = \supp(P)$ we have
$P(\Xi) = 1$. If the band holds at every point of $\Xi$, the violation set
is contained in $\Xi^{c}$, which has probability zero.

\textbf{(i) $\Rightarrow$ (ii).} Suppose the band holds $P$-a.e.\ but
fails at some $\bar\xi \in \Xi$. Two cases.

Upper bound violated: there is $i$ with $x_i > (1+\rho)\bar\xi_i$. Define
$g(\xi) := x_i - (1+\rho)\xi_i$, which is affine and therefore continuous
in $\xi$, and satisfies $g(\bar\xi) > 0$. By continuity there is
$\epsilon > 0$ with $g(\xi) > 0$ for all $\xi \in B_\epsilon(\bar\xi)$.
Because $\bar\xi \in \supp(P)$, every ball centred at $\bar\xi$ has
positive measure, so $P(B_\epsilon(\bar\xi)) > 0$. Every point of
$B_\epsilon(\bar\xi)$ violates the band, so the violation set has positive
probability, contradicting (i).

Lower bound violated: there is $i$ with $x_i < (1-\rho)\bar\xi_i$. Apply
the same argument to $\tilde g(\xi) := (1-\rho)\xi_i - x_i$, again affine
and continuous with $\tilde g(\bar\xi) > 0$.

Hence no such $\bar\xi$ exists and the band holds on all of $\Xi$.

\textbf{(ii) $\Leftrightarrow$ (iii).} Fix $i$. The requirement
$x_i \le (1+\rho)\xi_i$ for every $\xi \in \Xi$ is equivalent to
$x_i \le \inf_{\xi \in \Xi} (1+\rho)\xi_i$. Since $1+\rho > 0$ and $\Xi$
is a compact polytope, the infimum is attained and equals
$(1+\rho)\min_{\xi\in\Xi}\xi_i$. Likewise $x_i \ge (1-\rho)\xi_i$ for
every $\xi \in \Xi$ is equivalent to
$x_i \ge \sup_{\xi\in\Xi}(1-\rho)\xi_i = (1-\rho)\max_{\xi\in\Xi}\xi_i$,
where the direction of the inequality is preserved because
$1-\rho > 0$ for $\rho \in (0,1)$.
\end{proof}

Write $\underline{\xi}_i := \min_{\xi \in \Xi}\xi_i$ and
$\overline{\xi}_i := \max_{\xi \in \Xi}\xi_i$. Each is the optimal value of
a linear program over the polytope $\Xi$, so both are computable. The
constrained model becomes a finitely constrained problem:
\begin{align}
  \minimize_{x \ge 0} \quad & \mathbb{E}_{P}\big[ f(x,\xi) \big] \\
  \text{s.t.}\quad
    & (1-\rho)\,\overline{\xi}_i \;\le\; x_i \;\le\;
      (1+\rho)\,\underline{\xi}_i \qquad \forall\, i \in N .
  \label{eq:bandfinite}
\end{align}

\begin{remark}[When the band is empty]
\label{rem:feas}
Constraint \eqref{eq:bandfinite} is consistent for product $i$ if and only
if $(1-\rho)\overline{\xi}_i \le (1+\rho)\underline{\xi}_i$, that is
\begin{equation}
  \rho \;\ge\; \rho^{*}_i \;:=\;
  \frac{\overline{\xi}_i - \underline{\xi}_i}
       {\overline{\xi}_i + \underline{\xi}_i},
  \qquad\text{and the model is feasible iff }
  \rho \ge \max_{i \in N} \rho^{*}_i .
\end{equation}
This is the price of insisting on an almost-sure constraint with a
here-and-now decision: if the band is narrower than the spread of demand
it has to cover, no order quantity works at all. Section~1.7 confirms the
threshold numerically.
\end{remark}

\subsection*{1.6\quad(Part 3d) Comment}

No answer required. The point stands: an uncertain constraint required to
hold almost surely is a robust constraint in disguise, and the two usual
escapes are to relax it into a chance constraint, or to add second-stage
recourse so feasibility can be restored after $\xi$ is seen.

\subsection*{1.7\quad Classification and numerical check}

\begin{center}
\begin{tabular}{@{}ll@{}}
\toprule
Question (taxonomy, slides 8--10) & Answer for the SAA model \\
\midrule
Constrained?                & Yes \\
Variable type and size      & Continuous, $n + nm$ \\
Objective type              & Linear after the epigraph split;
                              piecewise linear convex before it \\
Constraint type and size    & Linear, $2nm$ rows \\
Parameters                  & Stochastic, distribution $P$ \\
Stages and players          & Single stage, one player \\
\bottomrule
\end{tabular}
\end{center}

A test instance with $n = 3$, independent uniform marginals on the box
$\prod_i [\underline{\xi}_i,\overline{\xi}_i]$, zero purchase cost and
\[
  h = (1.0,\,2.0,\,1.5), \qquad b = (4.0,\,3.0,\,5.0), \qquad
  \underline{\xi} = (10,5,8), \qquad \overline{\xi} = (30,25,20),
\]
admits a closed-form optimum. With $c=0$ the first-order condition for
product $i$ is $(h_i+b_i)F_i(x_i) = b_i$, so the critical fractile is
$F_i(x^{*}_i) = b_i/(b_i+h_i)$ and, for a uniform marginal,
\[
  x^{*}_i = \underline{\xi}_i
    + (\overline{\xi}_i - \underline{\xi}_i)\frac{b_i}{b_i+h_i},
  \qquad
  \mathbb{E}_P[f(x^{*},\xi)]
    = \sum_{i \in N}
      \frac{(\overline{\xi}_i-\underline{\xi}_i)\,h_i b_i}{2(h_i+b_i)} .
\]
This gives $x^{*} = (26.0,\,17.0,\,17.2308)$ and optimal expected cost
$26.923077$. The SAA converges to it:

\begin{center}
\begin{tabular}{@{}rrrr@{}}
\toprule
$m$ & SAA value (in sample) & $\mathbb{E}_P[f(\hat x_m,\xi)]$ &
  $\max_i |\hat x_{m,i} - x^{*}_i|$ \\
\midrule
50   & 24.2127 & 27.1680 & 0.9968 \\
200  & 27.2641 & 27.0311 & 0.9180 \\
1000 & 26.6339 & 26.9456 & 0.4228 \\
3000 & 26.5452 & 26.9271 & 0.1090 \\
\bottomrule
\end{tabular}
\end{center}

The downward bias of the in-sample column is visible at $m=50$, where the
SAA value $24.21$ sits below the true optimum $26.92$; the out-of-sample
column approaches the optimum from above, as it must.

For the band, the box support gives $\underline{\xi} = (10,5,8)$ and
$\overline{\xi} = (30,25,20)$, so
$\rho^{*} = (0.5000,\,0.6667,\,0.4286)$ and the binding threshold is
$\rho^{*}_2 = 2/3$. Solved at $m=1000$: the model is infeasible at
$\rho = 0.3,\ 0.45,\ 0.6$ and at $\rho = 0.6666$, and becomes feasible at
$\rho = 0.6668$ with value $51.8117$, against $26.5452$ unconstrained.
The threshold predicted by Remark~\ref{rem:feas} is reproduced exactly.

%=======================================================================
\newpage
\section*{Problem 2.\quad Robust Multidimensional Newsvendor
  \quad[40 pts]}
%=======================================================================

\subsection*{2.1\quad(Part 1) The robust model}

Replace the expectation by the worst case over the support:
\begin{align}
  \minimize_{x} \quad & \maximize_{\xi \in \Xi}\; f(x,\xi)
    \;=\; \maximize_{\xi \in \Xi} \sum_{i\in N}
      \Big( h_i(x_i-\xi_i)^{+} + b_i(\xi_i-x_i)^{+} \Big) \\
  \text{s.t.}\quad & x_i \ge 0 \qquad \forall\, i \in N .
\end{align}
No reformulation is performed. As written the inner problem has infinitely
many candidate $\xi$, so this is a semi-infinite program.

\subsection*{2.2\quad(Part 2) An equivalent LP with finitely many rows}

By Lemma~\ref{lem:convex}, $f(x,\cdot)$ is convex on the polytope $\Xi$. A
convex function maximised over a polytope attains its maximum at an
extreme point, so for every $x$
\begin{equation}
  \maximize_{\xi \in \Xi} f(x,\xi)
  \;=\; \max_{k \in \{1,\dots,K\}} f(x, v^{k}),
  \qquad \ext(\Xi) = \{v^1,\dots,v^K\}.
  \label{eq:vertexmax}
\end{equation}
Introduce $\tau \in \mathbb{R}$ for the worst-case cost and
$t^{k}_i \in \mathbb{R}$ for the inner maximum of product $i$ at vertex
$k$, then use \eqref{eq:costmax}:
\begin{align}
  \minimize_{x,\,\tau,\,t} \quad & \tau \\
  \text{s.t.}\quad
    & \tau \;\ge\; \sum_{i \in N} t^{k}_i
      && \forall\, k \in \{1,\dots,K\}
      \label{eq:p2epi}\\
    & t^{k}_i \;\ge\; h_i\,(x_i - v^{k}_i)
      && \forall\, i \in N,\; k \in \{1,\dots,K\} \\
    & t^{k}_i \;\ge\; b_i\,(v^{k}_i - x_i)
      && \forall\, i \in N,\; k \in \{1,\dots,K\} \\
    & x_i \ge 0 && \forall\, i \in N .
\end{align}
Size: $1 + n + nK$ variables and $K(2n+1)$ rows. Finite, and an LP.

\subsection*{2.3\quad(Part 3) When the extreme points cannot be enumerated}

$K$ is typically exponential in $n$, so \eqref{eq:p2epi} should be
generated on demand rather than written out. Use delayed constraint
generation.

\textbf{Master.} Keep a working set $W \subseteq \ext(\Xi)$ and solve the
LP of Section~2.2 restricted to $k \in W$. Dropping rows relaxes the
problem, so the master value is a lower bound on the robust optimum,
provided every element of $W$ is a genuine point of $\Xi$.

\textbf{Subproblem (separation).} Given the master solution
$(\hat x, \hat\tau)$, compute
$\;\mu(\hat x) := \max_{\xi \in \Xi} f(\hat x, \xi)$. This is an upper
bound on the robust optimum. If $\mu(\hat x) \le \hat\tau + \epsilon$,
stop: no extreme point is violated and $\hat x$ is optimal. Otherwise add
a maximiser to $W$ and re-solve.

\textbf{Solving the subproblem.} It maximises a convex function over a
polytope, which is NP-hard in general. Here the structure helps. Using
\eqref{eq:costmax} and the fact that the per-product choices are
independent,
\[
  \max_{\xi\in\Xi} \sum_{i\in N}
     \max\{ h_i(\hat x_i - \xi_i),\, b_i(\xi_i - \hat x_i) \}
  \;=\;
  \max_{\sigma \in \{H,B\}^{n}} \;\; \max_{\xi \in \Xi}
     \sum_{i \in N} \ell^{\sigma_i}_i(\xi),
\]
where $\ell^{H}_i(\xi) = h_i(\hat x_i-\xi_i)$ and
$\ell^{B}_i(\xi) = b_i(\xi_i-\hat x_i)$ are affine. For each fixed sign
pattern $\sigma$ the inner problem is a linear program over $\Xi$, whose
optimum is attained at an extreme point. So the separation problem is
$2^{n}$ linear programs, and the maximiser it returns is automatically an
extreme point and therefore a legitimate cut. Equivalently, it is a single
mixed-integer program with $n$ binaries,
\[
  \max \;\sum_i t_i \quad\text{s.t.}\quad
  t_i \le h_i(\hat x_i-\xi_i) + Mz_i,\;\;
  t_i \le b_i(\xi_i-\hat x_i) + M(1-z_i),\;\;
  \xi \in \Xi,\; z \in \{0,1\}^n,
\]
which is the better choice once $2^{n}$ becomes large.

\textbf{Termination.} $\ext(\Xi)$ is finite and each round adds a vertex
not already in $W$, because a vertex already in $W$ satisfies its own row
and cannot be strictly violated. So the loop stops after at most $K$
rounds, and in practice far sooner.

\begin{remark}[The seed must lie in $\Xi$]
\label{rem:seed}
Every element of $W$ must be a point of $\Xi$. Seeding with the
componentwise minimum $\underline{\xi} = (\underline{\xi}_1,\dots,
\underline{\xi}_n)$ is tempting and is wrong in general: each coordinate
minimum is attained in $\Xi$, but the vector collecting them need not lie
in $\Xi$ when the coordinates are coupled. Adding such a row makes the
master a restriction rather than a relaxation, and the procedure converges
to the wrong value. Section~2.9 reports an instance where this happens.
\end{remark}

\subsection*{2.4\quad(Part 4) The identity for $n=1$}

\begin{proposition}
For $h,b \ge 0$ and any $x,\xi \in \mathbb{R}$,
\[
  h(x-\xi)^{+} + b(\xi-x)^{+} \;=\; \max\{\, h(x-\xi),\; b(\xi-x) \,\}.
\]
\end{proposition}

\begin{proof}
This is Lemma~\ref{lem:identity} with $u = x-\xi$. Explicitly: if
$x \ge \xi$ then $(x-\xi)^{+} = x-\xi$ and $(\xi-x)^{+} = 0$, so the left
side is $h(x-\xi) \ge 0$, while $b(\xi-x) = -b(x-\xi) \le 0$, so the
maximum on the right is also $h(x-\xi)$. If $x < \xi$ then
$(x-\xi)^{+} = 0$ and $(\xi-x)^{+} = \xi-x$, so the left side is
$b(\xi-x) \ge 0$, while $h(x-\xi) \le 0$, so the maximum is $b(\xi-x)$.
The cases are exhaustive.
\end{proof}

The mechanism is that $x-\xi$ and $\xi-x$ have opposite signs, so with
$h,b \ge 0$ at most one of the two linear pieces is positive and the other
is nonpositive; adding the positive parts and taking the larger therefore
give the same number.

\subsection*{2.5\quad(Part 5) An alternative LP for $n=1$}

For $n=1$ a polytope $\Xi \subset \mathbb{R}_+$ is a closed interval
$[\underline{\xi},\overline{\xi}]$. Using the identity and exchanging the
two maxima,
\begin{align}
  \max_{\xi \in [\underline{\xi},\overline{\xi}]} f(x,\xi)
  &= \max_{\xi} \max\{ h(x-\xi),\, b(\xi-x) \} \notag\\
  &= \max\Big\{ \max_{\xi} h(x-\xi), \;\; \max_{\xi} b(\xi-x) \Big\}
  \;=\; \max\big\{ h(x-\underline{\xi}),\; b(\overline{\xi}-x) \big\},
  \label{eq:n1swap}
\end{align}
because $h(x-\xi)$ is nonincreasing in $\xi$ and so is largest at
$\xi = \underline{\xi}$, while $b(\xi-x)$ is nondecreasing in $\xi$ and so
is largest at $\xi = \overline{\xi}$. The robust problem is therefore
\begin{align}
  \minimize_{x \ge 0,\ \tau} \quad & \tau \\
  \text{s.t.}\quad & \tau \;\ge\; h\,(x - \underline{\xi}) \\
                   & \tau \;\ge\; b\,(\overline{\xi} - x) .
\end{align}
Two variables, two rows. Both rows are active at the optimum, so equating
them gives the closed form
\begin{equation}
  x^{*} = \frac{h\,\underline{\xi} + b\,\overline{\xi}}{h+b},
  \qquad
  \tau^{*} = \frac{h\,b\,(\overline{\xi}-\underline{\xi})}{h+b}.
  \label{eq:n1closed}
\end{equation}

The robust order is a convex combination of the two extreme demands,
weighted by the cost of the \emph{opposite} error: a large backlog cost
$b$ pulls $x^{*}$ towards $\overline{\xi}$, a large holding cost $h$ pulls
it towards $\underline{\xi}$. This is the robust counterpart of the
stochastic critical fractile $b/(b+h)$, and for a uniform marginal the two
coincide exactly, since
$\underline{\xi} + (\overline{\xi}-\underline{\xi})\,b/(b+h)
 = (h\underline{\xi} + b\overline{\xi})/(h+b)$.

\subsection*{2.6\quad(Part 6) Comparing the two reformulations}

\begin{center}
\begin{tabular}{@{}lrrl@{}}
\toprule
Formulation & Variables & Rows & Mechanism \\
\midrule
Extreme point (Section 2.2), $n=1$ &
  $1+n+nK = 4$ & $K(2n+1) = 6$ &
  enumerate vertices, then linearise at each \\
Identity based (Section 2.5) &
  $2$ & $2$ &
  resolve the inner maximum analytically \\
\bottomrule
\end{tabular}
\end{center}

The identity formulation is smaller, and the reason is structural rather
than incidental. The extreme-point route first enumerates the vertices and
only then linearises, paying one auxiliary variable $t^{k}_i$ and two rows
for every (product, vertex) pair. The identity route exchanges the two
maxima first, which lets each cost branch be sent to its own worst
demand analytically; after that no vertex enumeration and no auxiliary
variables are needed at all.

\textbf{What the measurement actually shows.} At $n=1$ both problems are
too small for the row count to matter. Measured on the instance of
Section~2.9: $6$ rows and $4$ columns solving in $0.0047$ s against $2$
rows and $2$ columns solving in $0.0042$ s, both returning
$\tau^{*} = 16.0$ at $x^{*} = 26.0$, which also matches
\eqref{eq:n1closed}. The difference is below timing noise, so the honest
claim for $n=1$ is that the identity formulation is structurally smaller,
not that it is measurably faster.

\textbf{Where the gap does bite.} The same exchange of maxima works for
any $n$ when $\Xi$ is a box $\prod_i[\underline{\xi}_i,\overline{\xi}_i]$,
because the inner maximisation then separates across products:
\begin{equation}
  \max_{\xi \in \Xi} f(x,\xi)
  = \sum_{i\in N} \max\big\{ h_i(x_i-\underline{\xi}_i),\;
      b_i(\overline{\xi}_i-x_i) \big\}
  \quad\Longrightarrow\quad
  \begin{array}{l}
    \min \sum_i \tau_i \\
    \tau_i \ge h_i(x_i-\underline{\xi}_i) \\
    \tau_i \ge b_i(\overline{\xi}_i-x_i)
  \end{array}
  \label{eq:boxsep}
\end{equation}
which is $2n$ rows, against $2^{n}$ vertices and $2^{n}(2n+1)$ rows for
the enumeration. Measured, with both formulations returning the same
optimal value at every size:

\begin{center}
\begin{tabular}{@{}rrrrrr@{}}
\toprule
$n$ & vertices $2^n$ & vertex rows & vertex solve (s) &
  identity rows & identity solve (s) \\
\midrule
2  & 4     & 20      & 0.004  & 4  & 0.004 \\
4  & 16    & 144     & 0.006  & 8  & 0.004 \\
6  & 64    & 832     & 0.024  & 12 & 0.005 \\
8  & 256   & 4\,352  & 0.161  & 16 & 0.006 \\
10 & 1\,024 & 21\,504 & 1.196  & 20 & 0.005 \\
12 & 4\,096 & 102\,400 & 47.328 & 24 & 0.005 \\
14 & 16\,384 & not built & --- & 28 & 0.006 \\
16 & 65\,536 & not built & --- & 32 & 0.005 \\
\bottomrule
\end{tabular}
\end{center}

Wall-clock timings are from one run on one machine and move by a few
per cent between runs; the row counts do not. The identity formulation
is flat in $n$ at this scale; the enumeration
doubles with every additional product and was abandoned past $n=12$, where
it already took $47.3$ s against $0.005$ s.

\textbf{The caveat that matters.} The separation \eqref{eq:boxsep} is
valid only when $\Xi$ is a box. For a general polytope the coordinates are
coupled and no single $\xi \in \Xi$ need attain every coordinate's own
worst extreme simultaneously, so \eqref{eq:boxsep} is then only an upper
bound and the extreme-point formulation, or constraint generation, is
required. Section~2.9 quantifies that gap on a coupled instance.

\subsection*{2.7\quad(Part 7) The robust model with the band constraint}

In a robust model the band is required over the whole uncertainty set by
definition, so no almost-sure argument is needed:
\begin{align}
  \minimize_{x \ge 0} \quad & \maximize_{\xi\in\Xi} f(x,\xi) \\
  \text{s.t.}\quad
    & (1-\rho)\,\xi_i \;\le\; x_i \;\le\; (1+\rho)\,\xi_i
      \qquad \forall\, i \in N,\; \forall\, \xi \in \Xi .
\end{align}

\subsection*{2.8\quad(Part 8) Yes, and here it is}

Both semi-infinite pieces reduce, for different reasons.

\emph{Objective.} $f(x,\cdot)$ is convex, so by \eqref{eq:vertexmax} the
worst case is attained at an extreme point and the objective needs only
the $K(2n+1)$ rows of Section~2.2.

\emph{Band.} For each $i$ the band is \emph{affine} in $\xi$, so its worst
case over $\Xi$ is attained at an extreme point as well, and by the same
computation as Proposition~\ref{prop:aerobust}(iii) it reduces to two
rows per product:
\[
  x_i \;\le\; (1+\rho)\,\underline{\xi}_i,
  \qquad
  x_i \;\ge\; (1-\rho)\,\overline{\xi}_i,
  \qquad
  \underline{\xi}_i = \min_{\xi\in\Xi}\xi_i, \;\;
  \overline{\xi}_i = \max_{\xi\in\Xi}\xi_i .
\]
Each bound is the value of one linear program over $\Xi$. The full finite
reformulation is
\begin{align}
  \minimize_{x,\tau,t} \quad & \tau \\
  \text{s.t.}\quad
    & \tau \ge \sum_{i\in N} t^{k}_i && \forall\, k \\
    & t^{k}_i \ge h_i(x_i - v^{k}_i),\quad
      t^{k}_i \ge b_i(v^{k}_i - x_i) && \forall\, i,\,k \\
    & (1-\rho)\,\overline{\xi}_i \;\le\; x_i \;\le\;
      (1+\rho)\,\underline{\xi}_i && \forall\, i \in N \\
    & x_i \ge 0 && \forall\, i \in N ,
\end{align}
with $K(2n+1) + 2n$ rows. The band therefore costs only $2n$ extra rows,
independent of $K$, and it is feasible under exactly the condition of
Remark~\ref{rem:feas}.

\subsection*{2.9\quad Classification and numerical results}

\begin{center}
\begin{tabular}{@{}ll@{}}
\toprule
Question (taxonomy, slides 8--10) & Answer for the extreme-point LP \\
\midrule
Constrained?             & Yes \\
Variable type and size   & Continuous, $1+n+nK$ \\
Objective type           & Linear \\
Constraint type and size & Linear, $K(2n+1)$ rows, $+2n$ with the band \\
Parameters               & Deterministic set $\Xi$; no distribution is used \\
Stages and players       & Single stage, two players: the decision maker
                           minimises, nature maximises over $\Xi$ \\
\bottomrule
\end{tabular}
\end{center}

Same $h, b$ as Problem 1, with a genuinely coupled support
\[
  \Xi = \Big\{ \xi :\; (10,5,8) \le \xi \le (30,25,20), \;\;
    45 \le \textstyle\sum_i \xi_i \le 60 \Big\},
\]
whose extreme points number $K = 12$. Results:

\begin{center}
\begin{tabular}{@{}lr@{}}
\toprule
Quantity & Value \\
\midrule
Robust optimal worst-case cost & $46.246154$ \\
Robust order quantity $x$ & $(30.0,\; 21.2,\; 19.2308)$ \\
LP size & $84$ rows ($=K(2n+1)$), $40$ columns \\
Separable box bound on the coordinate ranges & $53.846154$ \\
Value of the coupling & $7.600000$ \\
\bottomrule
\end{tabular}
\end{center}

The box bound overstates the worst case by $7.6$, because no single
$\xi\in\Xi$ can sit at every coordinate's own worst extreme at once. This
is the quantitative form of the caveat in Section~2.6.

\textbf{An observation worth recording.} If the lower total row is dropped
and only $\sum_i \xi_i \le 60$ is kept, the robust value equals the
separable box bound exactly ($53.846154$ both ways). The reason is that an
upper total row never binds at the robust optimum: the adversary's
cheapest retreat is the componentwise minimum $\underline{\xi}$, and if
$\Xi$ is nonempty then $\sum_i \underline{\xi}_i \le \sum_i \xi_i \le 60$,
so that retreat is always feasible. Only a \emph{lower} total row, which
forbids the retreat, makes the coupling bite. This also supplies the
counterexample promised in Remark~\ref{rem:seed}: here
$\underline{\xi} = (10,5,8)$ has total $23 < 45$ and is \emph{not} a point
of $\Xi$.

\textbf{Constraint generation.} Seeded at a genuine vertex of $\Xi$:

\begin{center}
\begin{tabular}{@{}rrrr@{}}
\toprule
Round & $|W|$ & Master (lower bound) & Separation (upper bound) \\
\midrule
1 & 1 & $0.000000$  & $135.000000$ \\
2 & 2 & $41.000000$ & $57.000000$ \\
3 & 3 & $43.707692$ & $54.707690$ \\
4 & 4 & $46.246154$ & $46.246155$ \\
\bottomrule
\end{tabular}
\end{center}

Four rounds, four of the twelve vertices ever added, and the same optimal
value as the full LP. Seeding instead at the componentwise minimum, which
is not in $\Xi$, converges to $53.107692$, a value \emph{above} the true
optimum, which is the signature of a master that has stopped being a
relaxation.

\textbf{The band.} With $\underline\xi = (10,5,8)$ and
$\overline\xi = (30,25,20)$ over $\Xi$, $\rho^{*} = (0.5,\,0.6667,\,
0.4286)$ and the threshold is $2/3$. Infeasible at $\rho = 0.4$ and at
$\rho = 0.6666$; feasible at $\rho = 0.6668$ with worst-case cost
$100.2515$, falling to $98.6154$ at $\rho = 0.70$, $91.5192$ at
$\rho=0.85$ and $85.0$ at $\rho = 1.0$, against $46.2462$ unconstrained.
The cost is decreasing in $\rho$, as it must be, since enlarging $\rho$
enlarges the feasible set.

\textbf{The identity, guarded numerically.} Over $200{,}000$ random
$(h,b,x,\xi)$ with $h,b \ge 0$, the largest discrepancy between the two
sides of Section~2.4 was $0$ to machine precision. The proof is the
case analysis; this only guards against a transcription slip.

%=======================================================================
\newpage
\section*{Problem 3.\quad CVaR Multidimensional Newsvendor
  \quad[15 pts, plus 5 extra]}
%=======================================================================

Throughout this problem $\Xi$ has finite support,
$\supp(P) = \{\xi^1,\dots,\xi^{|S|}\}$ with $P(\xi = \xi^s) = p_s > 0$.

\subsection*{3.1\quad(Part 1) The CVaR model}

For a loss $Z$ and level $\beta \in (0,1)$, write
$\VaR_\beta(Z) = \inf\{\eta : P(Z \le \eta) \ge \beta\}$ and let
$\CVaR_\beta(Z)$ be the mean of the worst $1-\beta$ tail of $Z$. The model,
with no reformulation performed, is
\begin{align}
  \minimize_{x} \quad & \CVaR_\beta\big[\, f(x,\xi) \,\big] \\
  \text{s.t.}\quad & x_i \ge 0 \qquad \forall\, i \in N .
\end{align}
Replacing $\mathbb{E}_P$ by $\CVaR_\beta$ changes what is being optimised,
not the decision structure: this is still a single-stage here-and-now
problem.

\subsection*{3.2\quad(Part 2) Epigraph reformulation and a
  master--subproblem scheme}

\textbf{Epigraph form.} By the Rockafellar--Uryasev representation,
\begin{equation}
  \CVaR_\beta(Z) \;=\; \min_{\eta \in \mathbb{R}}
    \Big\{\, \eta + \frac{1}{1-\beta}\,
      \mathbb{E}\big[(Z-\eta)^{+}\big] \Big\},
  \label{eq:ru}
\end{equation}
and the minimising $\eta$ is $\VaR_\beta(Z)$. Since the outer problem also
minimises, the two minimisations merge. Introducing
$u_s \ge 0$ for $(f(x,\xi^s)-\eta)^{+}$ and $t^{s}_i$ for the inner
maximum of \eqref{eq:costmax}:
\begin{align}
  \minimize_{x,\eta,u,t} \quad
    & \eta + \frac{1}{1-\beta} \sum_{s \in S} p_s\, u_s \\
  \text{s.t.}\quad
    & u_s \;\ge\; \sum_{i\in N} t^{s}_i \;-\; \eta && \forall\, s \in S \\
    & t^{s}_i \ge h_i(x_i - \xi^{s}_i),\quad
      t^{s}_i \ge b_i(\xi^{s}_i - x_i) && \forall\, i \in N,\; s \in S \\
    & u_s \ge 0,\quad x_i \ge 0 && \forall\, i \in N,\; s \in S .
\end{align}
This is a linear program with $1 + n + |S| + n|S|$ variables and
$|S|(2n+1)$ rows.

\textbf{Master--subproblem.} When $|S|$ is large, generate the scenario
rows on demand.

\emph{Master.} For a working set $W \subseteq S$, solve the LP above with
$s$ restricted to $W$. Dropping a scenario removes a nonnegative term
$p_s u_s/(1-\beta)$ from the objective, so the master value is a lower
bound.

\emph{Subproblem.} At the master solution $(\hat x, \hat\eta)$, evaluate
$f(\hat x, \xi^{s})$ for all $s \in S$ and form the true objective
\[
  \mathrm{UB} \;=\; \hat\eta + \frac{1}{1-\beta}
    \sum_{s \in S} p_s\,\big(f(\hat x,\xi^{s}) - \hat\eta\big)^{+},
\]
which upper-bounds the optimum by \eqref{eq:ru}. Stop when
$\mathrm{UB} - \mathrm{LB} \le \epsilon$; otherwise add to $W$ the
scenarios with the largest positive violation
$(f(\hat x,\xi^{s}) - \hat\eta)^{+}$ and re-solve. $S$ is finite and each
round adds at least one new index, so the scheme terminates.

\emph{Why it is efficient.} At the optimum $\hat\eta = \VaR_\beta$, so
every scenario with cost below $\hat\eta$ has $u_s = 0$ and contributes
nothing. Only the $(1-\beta)$ tail is active, and a working set of roughly
that size reproduces the full answer.

\begin{remark}[The master needs a lower bound on $\eta$]
\label{rem:unbdd}
Reducing $\eta$ by $\delta$ changes the master objective by
$-\delta\big(1 - \sum_{s\in W} p_s/(1-\beta)\big)$. So the master is
\emph{unbounded below} whenever the working set carries less than
$1-\beta$ of the probability mass, which is exactly the situation at the
start. Since $f \ge 0$ we have $\VaR_\beta(f) \ge 0$, so adding the valid
bound $\eta \ge 0$ repairs it without cutting off the optimum. This is not
an implementation detail: a naive scenario-dropping relaxation of
\eqref{eq:ru} has no finite value until $W$ reaches the tail mass.
\end{remark}

\subsection*{3.3\quad(Part 3, extra credit) Dualising the inner problem}

CVaR has the risk-envelope representation
\begin{equation}
  \CVaR_\beta(Z) \;=\; \max_{q \in \mathcal{Q}}\; \sum_{s\in S} q_s Z_s,
  \qquad
  \mathcal{Q} = \Big\{ q \ge 0 :\; \sum_{s\in S} q_s = 1,\;\;
    q_s \le \frac{p_s}{1-\beta} \Big\},
  \label{eq:envelope}
\end{equation}
so the CVaR problem is the two-level program
\[
  \minimize_{x \ge 0} \;\; \maximize_{q \in \mathcal{Q}} \;\;
    \sum_{s\in S} q_s\, f(x,\xi^{s}).
\]
For fixed $x$ the inner problem is a linear program in $q$. Attach the
free multiplier $\eta$ to $\sum_s q_s = 1$ and $\lambda_s \ge 0$ to
$q_s \le p_s/(1-\beta)$. Its dual is
\[
  \minimize_{\eta,\,\lambda \ge 0} \;\;
    \eta + \frac{1}{1-\beta}\sum_{s\in S} p_s \lambda_s
  \qquad\text{s.t.}\qquad
  \eta + \lambda_s \;\ge\; f(x,\xi^{s}) \quad \forall\, s \in S,
\]
the row for $s$ coming from the primal variable $q_s \ge 0$. The inner LP
is feasible and bounded, so strong duality holds and the maximum may be
replaced by this minimum. Since the outer problem also minimises, the two
merge into one level:
\begin{align}
  \minimize_{x,\eta,\lambda,t} \quad
    & \eta + \frac{1}{1-\beta}\sum_{s\in S} p_s \lambda_s \\
  \text{s.t.}\quad
    & \eta + \lambda_s \;\ge\; \sum_{i \in N} t^{s}_i
      && \forall\, s \in S \\
    & t^{s}_i \ge h_i(x_i - \xi^{s}_i),\quad
      t^{s}_i \ge b_i(\xi^{s}_i - x_i) && \forall\, i \in N,\; s \in S \\
    & \lambda_s \ge 0,\quad x_i \ge 0 && \forall\, i \in N,\; s \in S .
\end{align}

This is \emph{identical} to the epigraph model of Section~3.2 under
$\lambda_s = u_s$: the optimal $\lambda_s$ is
$(f(x,\xi^{s})-\eta)^{+}$. So the Rockafellar--Uryasev formulation is
precisely the linear-programming dual of the risk-envelope
representation, which is a satisfying way to see why \eqref{eq:ru} is
true rather than merely useful.

\subsection*{3.4\quad(Part 4) The CVaR model with the band constraint}

By Proposition~\ref{prop:aerobust} the band must hold on all of
$\supp(P)$. With finite support the argument is immediate and needs no
limit: each atom carries mass $p_s > 0$, so a violation at one atom is a
violation on a set of positive probability. Hence
\begin{align}
  \minimize_{x \ge 0} \quad & \CVaR_\beta\big[ f(x,\xi) \big] \\
  \text{s.t.}\quad
    & (1-\rho)\,\xi^{s}_i \le x_i \le (1+\rho)\,\xi^{s}_i
      && \forall\, i \in N,\; \forall\, s \in S ,
\end{align}
which collapses, exactly as before, to the $2n$ rows
\[
  (1-\rho)\max_{s\in S}\xi^{s}_i \;\le\; x_i \;\le\;
  (1+\rho)\min_{s\in S}\xi^{s}_i \qquad \forall\, i \in N,
\]
feasible under the condition of Remark~\ref{rem:feas} with the minimum and
maximum taken over the atoms. If instead one reads the uncertainty set as
$\conv\{\xi^{s}\}$, the bounds are unchanged, because a linear function
over a convex hull attains its extrema at the generating points.

\subsection*{3.5\quad Classification and numerical results}

\begin{center}
\begin{tabular}{@{}ll@{}}
\toprule
Question (taxonomy, slides 8--10) & Answer for the epigraph LP \\
\midrule
Constrained?             & Yes \\
Variable type and size   & Continuous, $1+n+|S|+n|S|$ \\
Objective type           & Linear \\
Constraint type and size & Linear, $|S|(2n+1)$ rows, $+2n$ with the band \\
Parameters               & Stochastic, finite support, known $p_s$ \\
Stages and players       & Single stage; one player, though
                           \eqref{eq:envelope} gives an adversarial
                           reading over $\mathcal{Q}$ \\
\bottomrule
\end{tabular}
\end{center}

Same $h,b$ and the same coupled $\Xi$ as Problem 2, with $|S| = 400$
equally likely scenarios drawn from $\Xi$ and $\beta = 0.95$.

\begin{center}
\begin{tabular}{@{}lr@{}}
\toprule
Quantity & Value \\
\midrule
Optimal $\CVaR_{0.95}$ & $38.647547$ \\
Order quantity $x$ & $(29.0513,\; 19.7828,\; 18.4501)$ \\
$\eta$ at the optimum, that is $\VaR_{0.95}$ & $36.721867$ \\
Scenarios strictly above $\eta$ & $20$ of $400$, as expected from
  $(1-\beta)|S| = 20$ \\
LP size & $2800$ rows \\
\bottomrule
\end{tabular}
\end{center}

\textbf{Ordering check.} Each model solved at its own optimum:
expectation $23.596948$, $\CVaR_{0.95}$ $38.647547$, worst case over the
scenario set $41.882681$. The required ordering
$\mathbb{E} \le \CVaR_\beta \le \max$ holds, which is the structural
sanity check for this family.

\textbf{Master--subproblem.} With $\eta \ge 0$ imposed per
Remark~\ref{rem:unbdd}:

\begin{center}
\begin{tabular}{@{}rrrr@{}}
\toprule
Round & $|W|$ & Master (lower bound) & True CVaR (upper bound) \\
\midrule
1 & 10 & $12.914440$ & $478.299416$ \\
2 & 20 & $32.404084$ & $181.320920$ \\
3 & 30 & $35.313817$ & $85.333800$ \\
4 & 40 & $38.617899$ & $38.680429$ \\
5 & 42 & $38.647547$ & $38.647547$ \\
\bottomrule
\end{tabular}
\end{center}

Five rounds and $42$ of $400$ scenarios, $10.5\%$, reproducing the full
LP value exactly. The jump between rounds 1 and 2 is the working set
crossing the tail mass of Remark~\ref{rem:unbdd}.

\textbf{The dualisation, checked.} At the optimal $x$, the risk-envelope
LP \eqref{eq:envelope} returns $38.647547$, equal to the CVaR computed
directly from the scenario costs. Its optimal $q$ has exactly $20$
nonzeros, all at the cap $p_s/(1-\beta) = 0.05$, which is the textbook
structure: CVaR puts the maximum admissible weight on the worst
$(1-\beta)$ fraction and zero elsewhere.

\textbf{The band.} Over the atoms, $\rho^{*}_i = (0.4978,\,0.6657,\,
0.4265)$, so the threshold is $0.6657$. Infeasible at $\rho = 0.4$ and
$\rho = 0.6647$; feasible at $\rho = 0.6667$ with CVaR $83.8378$, falling
to $79.2371$, $72.3338$ and $67.8643$ at $\rho = 0.75,\,0.90,\,1.00$,
against $38.6475$ unconstrained.

%=======================================================================
\newpage
\section*{Problem 4.\quad $\sigma$-Algebras and Pushforward Measures
  \quad[30 pts]}
%=======================================================================

\subsection*{4.1\quad(Part 1) The countable--cocountable $\sigma$-algebra}

\begin{proposition}
Let $\Omega$ be uncountable and
$\mathcal{F} := \{A \subseteq \Omega :\; A \text{ is countable, or }
\Omega\setminus A \text{ is countable}\}$. Then $\mathcal{F}$ is a
$\sigma$-algebra on $\Omega$.
\end{proposition}

\begin{proof}
\textbf{(i) $\emptyset \in \mathcal{F}$.} The empty set is finite, hence
countable, so $\emptyset \in \mathcal{F}$. (Also $\Omega \in \mathcal{F}$,
since $\Omega\setminus\Omega = \emptyset$ is countable.)

\textbf{(ii) Closed under complements.} Let $A \in \mathcal{F}$ and put
$A^{c} = \Omega\setminus A$. If $A$ is countable then
$\Omega\setminus A^{c} = A$ is countable, so $A^{c}$ has countable
complement and $A^{c} \in \mathcal{F}$. If instead $A^{c}$ is countable
then $A^{c} \in \mathcal{F}$ directly. Either way $A^{c} \in \mathcal{F}$.

\textbf{(iii) Closed under countable unions.} Let
$\{A_k\}_{k=1}^{\infty} \subseteq \mathcal{F}$ and $A = \bigcup_k A_k$.

\emph{Case 1: every $A_k$ is countable.} Then $A$ is a countable union of
countable sets, hence countable, so $A \in \mathcal{F}$.

\emph{Case 2: some $A_{k_0}$ has countable complement.} By De Morgan,
$\Omega\setminus A = \bigcap_k (\Omega\setminus A_k) \subseteq
\Omega\setminus A_{k_0}$. A subset of a countable set is countable, so
$\Omega\setminus A$ is countable and $A \in \mathcal{F}$.

The two cases are exhaustive, so $\mathcal{F}$ is closed under countable
unions and is a $\sigma$-algebra.
\end{proof}

\begin{remark}
Uncountability of $\Omega$ is what makes $\mathcal{F}$ interesting rather
than trivial: it guarantees that no set is both countable and
cocountable, so the two branches of the definition do not collapse, and
$\mathcal{F} \ne 2^{\Omega}$ whenever $\Omega$ has a subset that is
neither countable nor cocountable, as $[0,1]$ does.
\end{remark}

\subsection*{4.2\quad(Part 2) The $\sigma$-algebra generated by
  $\mathcal{F}_0$}

Let $\Omega = \{1,2,3,4\}$ and
$\mathcal{F}_0 = \big\{\{1\},\{3,4\}\big\}$. Since $\Omega$ is finite,
countable unions are finite unions and $\sigma(\mathcal{F}_0)$ is the
smallest family containing $\mathcal{F}_0$ that is closed under
complement and union.

Build it. We must have $\emptyset$ and $\Omega$. Complements force
$\{1\}^{c} = \{2,3,4\}$ and $\{3,4\}^{c} = \{1,2\}$. The union
$\{1\}\cup\{3,4\} = \{1,3,4\}$ must be present, hence so must its
complement $\{2\}$. The family now has the three disjoint sets
$\{1\},\{2\},\{3,4\}$, which partition $\Omega$; these are the
\emph{atoms}. Every union of atoms is forced, and the family of all such
unions is already closed under complement and union, so the construction
terminates:
\[
  \sigma(\mathcal{F}_0) \;=\;
  \Big\{\, \emptyset,\;
    \{1\},\; \{2\},\; \{3,4\},\;
    \{1,2\},\; \{1,3,4\},\; \{2,3,4\},\;
    \{1,2,3,4\} \,\Big\},
  \qquad |\sigma(\mathcal{F}_0)| = 8 .
\]

There are $3$ atoms and $2^{3} = 8$ members, as it must be for a
$\sigma$-algebra on a finite set. Note that $\{3\} \notin
\sigma(\mathcal{F}_0)$: the generators cannot distinguish the outcomes
$3$ and $4$, so no set in the generated $\sigma$-algebra separates them.
This is the standard reading of a $\sigma$-algebra as the information
available: $\sigma(\mathcal{F}_0)$ is exactly the collection of events
whose occurrence can be decided from knowing the answers to
``is the outcome in $\{1\}$?'' and ``is the outcome in $\{3,4\}$?''.

\subsection*{4.3\quad(Part 3) Closure under countable intersections}

\begin{proposition}
Let $\mathcal{F}$ be a $\sigma$-algebra on $\Omega$ and
$\{A_k\}_{k=1}^{\infty} \subseteq \mathcal{F}$. Then
$\bigcap_{k=1}^{\infty} A_k \in \mathcal{F}$.
\end{proposition}

\begin{proof}
By De Morgan,
\[
  \bigcap_{k=1}^{\infty} A_k
  \;=\; \Big( \bigcup_{k=1}^{\infty} A_k^{c} \Big)^{c}.
\]
Each $A_k \in \mathcal{F}$, so $A_k^{c} \in \mathcal{F}$ by closure under
complements. The countable union $\bigcup_k A_k^{c}$ is then in
$\mathcal{F}$ by closure under countable unions. Its complement is in
$\mathcal{F}$ by closure under complements again. Hence
$\bigcap_k A_k \in \mathcal{F}$.
\end{proof}

Only the three defining properties are used, and each is used exactly
where stated, which is why a $\sigma$-algebra is not required to list
intersections separately in its definition.

\subsection*{4.4\quad(Part 4) Pushforward of $\mathrm{Unif}[-1,1]$ under
  $X(\xi) = \xi^{2}$}

Let $\xi \sim P = \mathrm{Unif}[-1,1]$, so $P$ has density $\tfrac12$ on
$[-1,1]$, and let $X(\xi) = \xi^{2}$ with pushforward
$P_X = P \circ X^{-1}$.

\paragraph{(a) Support.} $X$ maps $[-1,1]$ onto $[0,1]$, and every
$t \in [0,1]$ has a preimage. For any $t \in [0,1]$ and $\epsilon>0$, the
set $X^{-1}\big((t-\epsilon,t+\epsilon)\big)$ contains a nondegenerate
interval of $\xi$ values and therefore has positive $P$-measure, so
$P_X\big((t-\epsilon,t+\epsilon)\big) > 0$. For $t \notin [0,1]$ a small
enough neighbourhood has preimage $\emptyset$. Hence
\[
  \supp(P_X) = [0,1].
\]

\paragraph{(b) Cumulative distribution function.} For $t \ge 0$,
\[
  X^{-1}\big((-\infty,t]\big)
  = \{\xi : \xi^{2} \le t\}
  = [-\sqrt{t},\,\sqrt{t}],
\]
so by the definition $P_X(A) = P(X^{-1}(A))$,
\[
  F_X(t) \;=\; P_X\big((-\infty,t]\big)
  \;=\; P\big([-\sqrt{t},\sqrt{t}\,] \cap [-1,1]\big)
  \;=\; \tfrac12 \cdot 2\sqrt{t} \;=\; \sqrt{t}
  \qquad \text{for } 0 \le t \le 1 .
\]
Collecting the cases,
\[
  F_X(t) =
  \begin{cases}
    0, & t < 0,\\[2pt]
    \sqrt{t}, & 0 \le t \le 1,\\[2pt]
    1, & t > 1 .
  \end{cases}
\]
$F_X$ is continuous and nondecreasing, with $F_X(0)=0$ and $F_X(1)=1$.

\paragraph{(c) Probability density function.} Differentiating on the
interior,
\[
  f_X(t) \;=\; \frac{\mathrm{d}}{\mathrm{d}t}\sqrt{t}
  \;=\; \frac{1}{2\sqrt{t}}, \qquad t \in (0,1),
\]
and $f_X(t) = 0$ otherwise. It integrates to one:
$\int_0^1 \tfrac{1}{2\sqrt t}\,\mathrm{d}t = \big[\sqrt t\,\big]_0^1 = 1$.
The density is unbounded as $t \downarrow 0$ but remains integrable, which
is the analytic counterpart of $F_X$ having infinite slope at the origin.
Squaring compresses the mass near $\pm 1$ and piles it up near $0$.

\subsection*{4.5\quad Numerical check}

The $\sigma$-algebra of Section~4.2 was also generated programmatically by
closing $\{\{1\},\{3,4\}\}$ under complement and union: the fixed point
has $8$ members, matching the list above exactly, with atoms
$\{1\},\{2\},\{3,4\}$, and it passes a direct test for closure under
complement and union. All $28$ pairwise intersections lie inside it,
consistent with Section~4.3.

For the pushforward, $4{,}000{,}000$ draws of $\xi \sim
\mathrm{Unif}[-1,1]$ give an empirical CDF of $X = \xi^{2}$ agreeing with
$\sqrt{t}$ to a worst deviation of $0.000364$ over a grid of $21$ points,
against a Monte Carlo standard error of about $0.0005$. The identity
$P(X \le t) = P(-\sqrt t \le \xi \le \sqrt t)$ holds exactly in the
sample. Moments confirm the density: $\mathbb{E}[X^{k}]$ came out
$0.333277$, $0.199984$ and $0.142872$ for $k=1,2,3$, against the exact
values $\int_0^1 t^{k}/(2\sqrt t)\,\mathrm{d}t = 1/(2k+1)$, that is
$1/3$, $1/5$ and $1/7$.

%=======================================================================
\newpage
\section*{EP3.\quad What are coherent risk measures?}
%=======================================================================

A risk measure assigns to a random loss $Z$ a single number
$\varrho(Z)$ saying how much that loss should worry us. Variance was the
classical answer and is a poor one for losses: it penalises favourable
deviations as hard as unfavourable ones, and it is not expressed in the
units of the loss. Artzner, Delbaen, Eber and Heath (1999) asked instead
which properties a sensible measure of loss \emph{ought} to have, and
called a measure \emph{coherent} when it has all four.

\textbf{The four axioms.} For losses $Z, Z_1, Z_2$ and constants
$c \in \mathbb{R}$, $\lambda \ge 0$:
\begin{enumerate}
  \item \emph{Monotonicity.} $Z_1 \le Z_2$ pointwise implies
    $\varrho(Z_1) \le \varrho(Z_2)$. A position that loses more in every
    state is riskier.
  \item \emph{Translation equivariance.} $\varrho(Z+c) = \varrho(Z)+c$.
    Adding a certain loss of $c$ raises the requirement by exactly $c$,
    so $\varrho$ is denominated in money.
  \item \emph{Positive homogeneity.} $\varrho(\lambda Z) =
    \lambda\varrho(Z)$. Doubling a position doubles its risk.
  \item \emph{Subadditivity.} $\varrho(Z_1+Z_2) \le \varrho(Z_1) +
    \varrho(Z_2)$. Merging two positions cannot create risk; this is the
    formal statement that diversification should never be punished.
\end{enumerate}
Axioms 3 and 4 together give convexity,
$\varrho(\lambda Z_1 + (1-\lambda)Z_2) \le \lambda\varrho(Z_1) +
(1-\lambda)\varrho(Z_2)$ for $\lambda \in [0,1]$, which is what makes a
coherent measure usable as an objective in a convex program.

\textbf{Why VaR fails and CVaR does not.} $\VaR_\beta$ satisfies axioms 1
to 3 but violates subadditivity. The standard counterexample is two
independent defaultable bonds, each with a small default probability
below $1-\beta$: separately each has $\VaR_\beta = 0$, but the
portfolio defaults with probability large enough to push $\VaR_\beta$
above zero, so $\varrho(Z_1+Z_2) > \varrho(Z_1)+\varrho(Z_2)$. VaR
therefore reports that splitting a portfolio in two reduces its risk,
which is the opposite of what diversification means. $\CVaR_\beta$ repairs
this: it averages the whole tail rather than reading off a single
quantile, and it satisfies all four axioms.

\textbf{The representation theorem, and why it matters here.} Every
coherent risk measure on a finite probability space is the support
function of a closed convex set of probability measures:
\[
  \varrho(Z) \;=\; \max_{q \in \mathcal{Q}} \; \mathbb{E}_{q}[Z]
\]
for some \emph{risk envelope} $\mathcal{Q}$. Coherent risk minimisation is
therefore robust optimisation in disguise, with $\mathcal{Q}$ as the
ambiguity set. The two extremes bracket everything in between: taking
$\mathcal{Q} = \{P\}$ recovers the expectation, and taking $\mathcal{Q}$
to be all distributions supported on $\Xi$ recovers the worst case. CVaR
sits between them with
\[
  \mathcal{Q}_{\CVaR} = \Big\{ q \ge 0 :\;
    \textstyle\sum_s q_s = 1,\;\; q_s \le p_s/(1-\beta) \Big\},
\]
which is the set used in Problem 3. That is exactly why the three models
in this homework line up as
$\mathbb{E} \le \CVaR_\beta \le \max$, confirmed numerically there as
$23.5969 \le 38.6475 \le 41.8827$: they are the same optimisation problem
over three nested ambiguity sets. It is also why dualising the inner
maximisation in Section~3.3 returns the Rockafellar--Uryasev epigraph
model: the dual of a support function is the epigraph form of the measure.

\textbf{What the axioms cost.} Positive homogeneity is the one most often
given up. It rules out any extra penalty for sheer size, yet a position
ten times larger may be harder to liquidate than ten separate positions.
Dropping it while keeping monotonicity, translation equivariance and
convexity gives the larger class of \emph{convex} risk measures, of which
the entropic risk measure $\tfrac1\gamma\log\mathbb{E}[e^{\gamma Z}]$ is
the standard example. For optimisation purposes convexity is the property
that actually matters, since it is what keeps the problem tractable;
coherence is the stronger requirement one imposes when the number must
also behave like a capital requirement.

\textbf{References.} P.~Artzner, F.~Delbaen, J.-M. Eber and D.~Heath,
``Coherent Measures of Risk'', \emph{Mathematical Finance} 9(3),
203--228, 1999. R.~T. Rockafellar and S.~Uryasev, ``Optimization of
Conditional Value-at-Risk'', \emph{Journal of Risk} 2(3), 21--41, 2000.

%=======================================================================
\vspace{1em}
\hrule
\vspace{0.6em}
{\small All numerical values quoted above were produced by the
accompanying Python files \texttt{p1\_stochastic.py},
\texttt{p2\_robust.py}, \texttt{p3\_cvar.py} and
\texttt{p4\_measure.py}, solved with PuLP and CBC. Every instance outside
the assignment statement was constructed for this report; the figures are
specific to those instances and are not claims about the general case.}

\end{document}
